\documentclass[12pt]{article}
\usepackage[utf8]{inputenc}
\usepackage[T1]{fontenc}
\usepackage[margin=1in]{geometry}
\usepackage{setspace}
\usepackage{xcolor}
\usepackage[colorlinks=true, urlcolor=blue!60!black, linkcolor=blue!60!black]{hyperref}
\setlength{\parindent}{2em}
\setlength{\parskip}{1.25em}
\renewcommand{\baselinestretch}{1.0}
\begin{document}
\begin{center}
{\LARGE \textit{City of Memphis v.\ Greene}, 451 U.S.\ 100 (1981)\par}
\vspace{0.8em}
{\large Case Brief --- Four-Page Edition\par}
\end{center}
\vspace{1em}

\section{Memory Jogger}

Memphis's closure of a street between an all-white neighborhood and a predominantly Black one, to reduce through traffic, did not violate 42 U.S.C.\ \S~1982 or the Thirteenth Amendment: it caused only a routine commuting inconvenience, without proof of discriminatory intent or impairment of Black property rights (\href{https://supreme.justia.com/cases/federal/us/451/100/}{Justia}).

\section{Detailed Case Facts}

Hein Park, a Memphis subdivision developed before World War II as an exclusive white neighborhood, was entirely white-owned; the area just north was predominantly Black. West Drive, a half-mile, two-lane street, ran through Hein Park from Jackson Avenue south toward Overton Park. In 1970 the city denied residents' request to close four streets, but traffic engineers noted that closing West Drive at Jackson Avenue would eliminate much through traffic. In 1973 the Hein Park Civic Association applied to close the northernmost 25 feet of West Drive and deed the strip to the abutting owners, citing through traffic, the safety of children walking to school, and ``traffic pollution.'' Despite petitions with roughly 1,000 signatures protesting that the closure symbolized a white neighborhood shutting its door on Black neighbors, the City Council approved it. Memphis had a documented history of resisting desegregation.

\section{Procedural History}

Black residents north of Jackson Avenue filed a class action under \S\S~1982 and 1983 and the Thirteenth and Fourteenth Amendments. The district court dismissed; the Sixth Circuit reversed. After trial, Judge McRae found a disproportionate impact on ``overwhelmingly black'' traffic but no denial of benefits, no discriminatory intent under \textit{Arlington Heights}, and no significant procedural departure. The Sixth Circuit again reversed, holding the closure a ``badge of slavery'' impairing Black property rights. The Supreme Court granted certiorari.

\section{Judicial Votes}

Reversed, \textbf{6--3} (decided April 20, 1981).
\begin{itemize}
\item \textbf{Majority:} Justice Stevens, joined by Chief Justice Burger and Justices Stewart, Powell, and Rehnquist.
\item \textbf{Concurring in the judgment:} Justice White.
\item \textbf{Dissent:} Justice Marshall, joined by Justices Brennan and Blackmun.
\end{itemize}

\section{Holding}

A facially neutral street closure serving legitimate traffic-safety and residential-tranquility interests, without proof of discriminatory intent, property depreciation, severe restriction of access to homes, or unequal denial of municipal benefits, does not impair property rights protected by \S~1982 and is not a ``badge or incident of slavery'' under the Thirteenth Amendment, even if it disparately inconveniences Black motorists. The Court left open whether \S~1982 requires purposeful discrimination and whether Section~1 of the Thirteenth Amendment, of its own force, reaches badges and incidents of slavery.

\section{Analysis of Opinions}

\textbf{Majority (Stevens).} The record did not support the Sixth Circuit's premises: southbound drivers merely had to turn a few blocks earlier; the barrier was a curb with a gap for municipal vehicles; one ``undesirable traffic'' remark did not show the Council's animus; and plaintiffs' own broker testified that values would not decline. Absent intent, \textit{Washington v.\ Davis} and \textit{Arlington Heights} foreclosed an equal protection claim. Section 1982 protects the equal right to acquire, hold, and use property; because Memphis neither denied Black neighborhoods comparable closures, depreciated their property, nor severely restricted access to their homes, requiring use of another street impaired no \S~1982 interest. Nor was the closure a badge of slavery: local governments may balance convenience against safety and tranquility, and ordinary traffic rules inevitably have disparate effects in racially concentrated neighborhoods.

\textbf{White (concurring in the judgment).} The Court should not have reweighed the record but answered the question presented. The 1866 Act targeted the Black Codes and intentional oppression, so \S~1982 requires purposeful discrimination; given the finding of no intent, reversal followed.

\textbf{Marshall (dissenting).} The barrier, placed exactly at the line between an all-white enclave and a Black community, was a ``monument to racial hostility.'' Marshall pointed to \textit{Arlington Heights} evidence the majority ignored: the ``undesirable traffic'' testimony, procedural irregularities, the city's never having closed a through street for traffic control, and its history of segregation. Known, massively disproportionate harm to Black property enjoyment demanded a justification Memphis did not supply, and the barrier also violated the Thirteenth Amendment.

\section{Examples: Future Applications}

\begin{enumerate}
\item \textbf{Same-Side: Traffic Calming.} A city makes a street through a white historic district one-way with speed humps to curb cut-through traffic from a Hispanic neighborhood with access to parallel arterials. With no animus or property impairment, it is a routine burden.
\item \textbf{Same-Side: Moved Bus Hub.} Relocating a transit hub adds four minutes to Black residents' commutes; indirect inconvenience without a property nexus fails.
\item \textbf{Opposite-Side: Locked Gate.} A town gates the only paved road to a Black enclave but gives a key to the lone white landowner. Severely restricting access to Black-owned homes impairs \S~1982 rights, as the Court's distinction of \textit{Jennings v.\ Patterson} indicates.
\item \textbf{Opposite-Side: Selective Closures.} A city grants cul-de-sacs to five white subdivisions but denies identical petitions from Black neighborhoods. Unequal municipal benefits violate \S~1982.
\item \textbf{Fence-Sitter: Sound Wall.} A city builds a wall with no pedestrian openings between a white subdivision and a Black apartment complex, citing noise but also complaints about ``foot traffic and crime.'' Pretext evidence is stronger than in \textit{Greene}, but the city will cite noise abatement and remaining vehicle access.
\end{enumerate}

\section{Critique}

\textbf{Doctrinal limbo.} As Justice White charged, by ducking the intent question and reweighing evidence, the Court left \S~1982 intent doctrine and Section~1 of the Thirteenth Amendment unresolved, just thirteen years after \textit{Jones v.\ Alfred H.\ Mayer Co.} \textbf{Antisubordination.} Critics say the majority's formalism echoed \textit{Plessy}'s dismissal of stigma, erased the history that created all-white Hein Park, and froze Section~1 litigation, which lower courts then confined to literal coerced labor.

\textbf{Originalism.} Textualists applaud the refusal to call a traffic barrier ``slavery'' and White's reading of the 1866 Act; some Reconstruction originalists argue Section~1 itself abolished racially exclusionary customs.

\section{Key Quotations}

\begin{quote}
``We merely hold that the impact of the closing of West Drive on nonresidents of Hein Park is a routine burden of citizenship; it does not reflect a violation of the Thirteenth Amendment.''
\end{quote}
\noindent --- Justice Stevens, for the Court, 451 U.S.\ at 129.

\begin{quote}
``Because urban neighborhoods are so frequently characterized by a common ethnic or racial heritage, a regulation's adverse impact on a particular neighborhood will often have a disparate effect on an identifiable ethnic or racial group. To regard an inevitable consequence of that kind as a form of stigma so severe as to violate the Thirteenth Amendment would trivialize the great purpose of that charter of freedom.''
\end{quote}
\noindent --- Justice Stevens, for the Court, 451 U.S.\ at 128.

\begin{quote}
``Petitioner city of Memphis, acting at the behest of white property owners, has closed the main thoroughfare between an all-white enclave and a predominantly Negro area of the city. The stated explanation for the closing is of a sort all too familiar: `protecting the safety and tranquility of a residential neighborhood' by preventing `undesirable traffic' from entering it. Too often in our Nation's history, statements such as these have been little more than code phrases for racial discrimination.''
\end{quote}
\noindent --- Justice Marshall, dissenting, 451 U.S.\ at 135--36.

\end{document}
